\begin{table}[htbp]
\centering
\small
\renewcommand{\arraystretch}{1.5}
\begin{tabular}{>{\raggedright\arraybackslash}p{5.5cm} >{\raggedright\arraybackslash}p{4.2cm} >{\raggedright\arraybackslash}p{4.5cm}}
\hline
\textbf{Requisito funzionale} & \textbf{Schema E/R} & \textbf{Schema relazionale} \\
\hline

Gestire l'anagrafica degli studenti &
Entità \emph{Studente} &
\texttt{studente} \\

Gestire il corpo docente &
Entità \emph{Docente} &
\texttt{docente} \\

Gestire i dipartimenti e i corsi di laurea &
Entità \emph{Dipartimento}, \emph{Corso di Laurea}; relazioni di afferenza e iscrizione &
\texttt{dipartimento}, \texttt{corso\_di\_laurea}, FK in \texttt{docente} e \texttt{studente} \\

Gestire gli insegnamenti &
Entità \emph{Insegnamento} &
\texttt{insegnamento} \\

Gestire le edizioni annuali degli insegnamenti &
Entità \emph{Edizione} &
\texttt{edizione} \\

Associare i docenti alle edizioni degli insegnamenti &
Relazione \emph{tiene} &
\texttt{insegnamento\_edizione} \\

Gestire il calendario delle lezioni &
Entità \emph{Lezione}; relazione \emph{si compone di} &
\texttt{lezione} \\

Gestire il piano di studio degli studenti &
Relazione \emph{nel piano di studio} &
\texttt{piano\_di\_studio} \\

Registrare gli esami sostenuti &
Entità \emph{Esame}; relazione \emph{sostiene} &
\texttt{esame} \\

Gestire le propedeuticità tra insegnamenti &
Relazione \emph{prerequisito} &
\texttt{prerequisito} \\

Gestire le abilitazioni dei docenti agli insegnamenti &
Relazione \emph{può insegnare} &
\texttt{abilitazione\_\allowbreak docente\_\allowbreak insegnamento} \\
\hline

\end{tabular}
\caption{Tracciabilità dei requisiti tra analisi, progettazione concettuale e progettazione logica.}
\label{tab:tracciabilita}
\end{table}